\section{Introduction}
A predicted probability vector and an uncertain compositional response are different statistical objects. A classifier may output a vector of probabilities while its observed outcome is a single label. Here the response itself is observed as a nonnegative vector summing to one. Examples include measured proportions and repeated-measurement distributions. Our coverage event is simultaneous inclusion of this entire vector in a region of the simplex; covering a label or each coordinate separately answers a different question.

Split conformal prediction provides a finite-sample marginal coverage guarantee under exchangeability, for a score fixed independently of calibration data. This leaves several choices open. Raw-coordinate and log-ratio distances produce different regions. A global calibration threshold averages over heterogeneous populations. Smaller regions can reflect undercoverage, and radii in different geometries have incompatible units. Finally, a guarantee averaging over calibration samples need not describe the performance of every realized calibration sample.

We study these choices in a controlled, synthetic compositional regression benchmark. The design uses a shared fitted predictor, independent repetitions of the full training--calibration--test experiment, and explicit subgroup and distribution-shift diagnostics. It separates four questions: whether coordinatewise calibration delivers simultaneous coverage; whether marginal coverage conceals subgroup failures; how geometry changes efficiency under a common reference measure; and which conclusions survive a controlled mixture shift, finite calibration ranks, and zero replacement. The contribution is an evaluation design and its reproducible evidence, rather than a new conformal guarantee. The controlled setting exposes mechanisms; it cannot establish performance for a particular application or a universal ranking of methods.

Our theory-based expectations are that correctly ranked joint scores achieve iid marginal coverage, naive coordinate calibration can miss the joint target, and heterogeneity can hide group undercoverage that becomes more consequential under mixture shift. We test whether scaling or group calibration improves this behavior and changes common-measure area; no geometry is presumed universally best. These hypotheses concern the controlled mechanisms, not all compositions.
